\chapter{Costos y Factibilidad Económica}

\section{3.3.3. Factibilidad Económica}
Dado que se trata de un proyecto académico y sin fines de lucro, se opta por una estrategia de desarrollo de bajo costo. 
Se emplean herramientas y plataformas de uso libre o con planes académicos gratuitos, como GitHub para el control de versiones, MongoDB en su modalidad gratuita, y Google Drive para el almacenamiento de evidencias. 
El desarrollo se realiza de forma local en los equipos personales de los integrantes, distribuyendo las tareas según los módulos definidos en el WBS. 
Esta estrategia permite garantizar la continuidad del proyecto sin generar gastos adicionales a la institución.

\section{Costos de desarrollo}
El desarrollo se ejecuta dentro del marco del curso \textit{Ingeniería en Sistemas II} de la Universidad Nacional, 
por lo que no existen costos asociados a contratación o licencias de software. 
Todas las herramientas utilizadas son gratuitas y de libre acceso, y el despliegue de las pruebas se realiza en entornos locales. 
De requerirse demostraciones institucionales, estas se efectuarán mediante el uso de cuentas académicas o servidores internos de la UNA.

\section{Costos de capital humano}
Las actividades técnicas están a cargo del equipo de desarrollo estudiantil, compuesto por:

\begin{itemize}
  \item \textbf{Francisco Mora Cabezas:} Desarrollo del backend del módulo \textit{Tiempos de Jornada}, diseño y definición del esquema correspondiente, integración visual y elaboración de la documentación general del proyecto.
  \item \textbf{Ángel Segura Méndez:} Desarrollo del backend del módulo \textit{Evidencias SINAES}, modelado de su base de datos y elaboración de documentación técnica específica del módulo.
\end{itemize}

Cada integrante es responsable del ciclo completo de su módulo, incluyendo la lógica del backend, pruebas unitarias y coordinación para la integración futura entre ambos sistemas. 
No se contemplan remuneraciones, dado que el trabajo forma parte del proceso académico dentro del curso. 
La participación docente se limita a la supervisión, seguimiento metodológico y validación de los avances.


\section{Cargos administrativos}
El sistema está diseñado para operar completamente en línea o en entornos de prueba locales, eliminando la necesidad de documentación física. 
Por ello, los costos administrativos son prácticamente nulos. 
Se prevé únicamente el uso de materiales de apoyo o presentaciones digitales dentro del entorno académico.

\section{Costos totales del proyecto}
El proyecto puede desarrollarse sin generar costos monetarios directos. 
No obstante, se estima un gasto potencial máximo de ₡25\,000 a ₡30\,000 (aproximadamente 40–50 USD) en caso de requerir:
\begin{itemize}
  \item Compra de un dominio institucional para futuras pruebas públicas.
  \item Certificados SSL u otros servicios asociados al despliegue de pruebas.
\end{itemize}
Estos costos son opcionales y no forman parte de la presente etapa del proyecto.

\section{Resumen consolidado de costos}

A continuación se presenta una tabla consolidada que resume todos los costos asociados al desarrollo e implementación del Sistema de Gestión Integral Universitario:

\begin{table}[H]
\centering
\begin{tabular}{|l|r|p{5cm}|}
\hline
\textbf{Categoría de Costo} & \textbf{Monto (₡)} & \textbf{Observaciones} \\
\hline
\hline
\multicolumn{3}{|l|}{\textbf{I. Costos de Desarrollo}} \\
\hline
Herramientas y software & 0.00 & Uso de tecnologías open source \\
Plataformas (GitHub, MongoDB Free, Vercel) & 0.00 & Planes académicos gratuitos \\
IDE y entornos de desarrollo & 0.00 & VS Code, Node.js (gratuitos) \\
\hline
\textbf{Subtotal Desarrollo} & \textbf{0.00} & \\
\hline
\hline
\multicolumn{3}{|l|}{\textbf{II. Costos de Capital Humano}} \\
\hline
Desarrollo técnico (576 horas) & 0.00 & Proyecto académico \\
Supervisión docente & 0.00 & Dentro del curso ISII \\
Reuniones y coordinación & 0.00 & Herramientas virtuales gratuitas \\
\hline
\textbf{Subtotal Capital Humano} & \textbf{0.00} & \\
\hline
\hline
\multicolumn{3}{|l|}{\textbf{III. Costos Administrativos}} \\
\hline
Documentación física & 0.00 & Todo en formato digital \\
Comunicaciones & 0.00 & Google Meet, Discord, WhatsApp \\
Almacenamiento (Google Drive) & 0.00 & Cuentas académicas \\
\hline
\textbf{Subtotal Administrativo} & \textbf{0.00} & \\
\hline
\hline
\multicolumn{3}{|l|}{\textbf{IV. Costos Opcionales (Futuros)}} \\
\hline
Dominio institucional (.cr) & 10,000 - 15,000 & Si se requiere despliegue público \\
Certificado SSL & 5,000 - 10,000 & Si se requiere HTTPS personalizado \\
Hosting demo (opcional) & 0 - 10,000 & Alternativa: usar Vercel gratuito \\
\hline
\textbf{Subtotal Opcional} & \textbf{15,000 - 35,000} & \\
\hline
\hline
\textbf{TOTAL GARANTIZADO} & \textbf{₡0.00} & \textbf{Costo real del proyecto} \\
\textbf{TOTAL MÁXIMO ESTIMADO} & \textbf{₡35,000} & \textbf{Solo si se requieren servicios adicionales} \\
\hline
\end{tabular}
\caption{Consolidado General de Costos del Proyecto}
\label{tab:costos_consolidado}
\end{table}

\textbf{Nota importante sobre costo-beneficio:} 

Aunque el costo directo del proyecto es de ₡0.00, el valor generado es significativo. La Universidad Nacional actualmente invierte aproximadamente ₡58 millones anuales en gestión manual de evidencias SINAES (según estimaciones institucionales de horas-persona dedicadas a búsqueda, organización y auditoría de documentos). 

Adicionalmente, el módulo de Tiempos de Jornada optimiza un proceso que actualmente consume entre 40-60 horas por ciclo académico en cálculos manuales, revisiones y correcciones de asignaciones docentes. La automatización de este proceso representa un ahorro estimado de ₡2.4 millones anuales en productividad administrativa.

Por lo tanto, la relación costo-beneficio es extraordinariamente favorable: **inversión de ₡0.00 con retorno anual estimado de ₡60+ millones en ahorro de tiempo y optimización de recursos humanos**.

\section{Conclusiones de la factibilidad económica}
El proyecto es económicamente viable y presenta una relación costo-beneficio excepcional. 
Aprovecha herramientas gratuitas de código abierto, no requiere inversión inicial y mantiene un costo de mantenimiento prácticamente nulo. 
Al estar alineado con los principios de software libre y educación abierta, contribuye al fortalecimiento institucional sin demandar recursos financieros adicionales. 

El único costo potencial (₡15,000-₡35,000) surgiría únicamente si la institución decidiera en el futuro implementar un despliegue público con dominio personalizado, certificado SSL profesional y servicios premium de hosting. Sin embargo, este gasto sería completamente opcional, ya que el sistema puede operar indefinidamente con las herramientas gratuitas actuales.

En términos de costo-beneficio, la relación es altamente favorable: el impacto esperado —automatización, transparencia, eficiencia administrativa y ahorro de ₡60+ millones anuales— se logra con una inversión real de cero colones.
+++